\rtmtight
\begin{tblr}{width=\textwidth,colspec={|Q[c,m,2.1cm]|X[l,m]|},hlines,vlines,hspan=minimal,rowsep=1.5pt,colsep=3pt,row{1}={bg=headgray,font=\bfseries}}
\SetCell{c,m}\hfil\logo\hfil & \textbf{POLITEKNIK NEGERI MALANG}\par\textbf{JURUSAN TEKNOLOGI INFORMASI}\par\textbf{PROGRAM STUDI: D4 TEKNIK INFORMATIKA} \\
\SetCell[c=2]{l} \textbf{RENCANA TUGAS MAHASISWA (RTM) DAN RUBRIK PENILAIAN} & \\
Nama Mata Kuliah & Pemrograman Web Lanjut \\
Kode Mata Kuliah & RTI254007 \quad \textbf{sks / Jam perminggu:} 3 SKS/6 jam \quad \textbf{Semester:} 4 \\
Dosen Pengampu & \begin{enumerate}\item Farid Angga Pribadi, S.Kom., M.Kom.\item Moch. Zawaruddin Abdullah, S.ST., M.Kom.\item Dimas Wahyu Wibowo, S.T., M.T.\item Habibie Ed Dien, S.Kom., M.T.\item Endah Septa Sintiya, S.Pd., M.Kom.\item Muhammad Unggul Pamenang, S.St., M.T.\end{enumerate} \\
\end{tblr}
\assessmentblock{Tugas Instalasi, Routing, Migration, dan Blade}{Hands-on Practice}{SCPMK0209-03401, SCPMK0603-03402}{Mahasiswa menyelesaikan tugas instalasi, routing, migration, dan blade sesuai Sub-CPMK dan konteks mata kuliah.}{Menganalisis kebutuhan, mengerjakan artefak praktik/proyek, menguji hasil, mendokumentasikan evidence, dan mempresentasikan keputusan bila diminta.}{Laporan, artefak digital, repository/prototype, evidence pengujian, dan/atau bahan presentasi.}{Kesiapan proyek Laravel dan konfigurasi & 14 & Instalasi Laravel, environment, database connection, package, dan struktur proyek siap digunakan. & Proyek berjalan, tetapi konfigurasi environment atau struktur awal belum lengkap. & Proyek tidak berjalan atau konfigurasi dasar salah. \\ \\
Routing, controller, migration, dan model & 12.3 & Route, controller, migration, model, dan relasi dasar sesuai kebutuhan modul serta konsisten. & Komponen utama tersedia, tetapi relasi, migration, atau controller belum sepenuhnya konsisten. & Routing/model/migration tidak terhubung atau tidak sesuai fitur. \\ \\
Blade, validasi, dan evidence praktik & 8.7 & Blade view, form, validasi, feedback error, dan screenshot/commit praktik terdokumentasi jelas. & View dan form berjalan, tetapi validasi atau evidence praktik masih terbatas. & View tidak mendukung alur fitur atau tidak ada bukti praktik. \\}

\assessmentblock{UTS Modul CRUD Laravel}{UTS praktik}{SCPMK0209-03401, SCPMK0603-03402}{Mahasiswa menyelesaikan uts modul crud laravel sesuai Sub-CPMK dan konteks mata kuliah.}{Menganalisis kebutuhan, mengerjakan artefak praktik/proyek, menguji hasil, mendokumentasikan evidence, dan mempresentasikan keputusan bila diminta.}{Laporan, artefak digital, repository/prototype, evidence pengujian, dan/atau bahan presentasi.}{Kelengkapan CRUD dan validasi Laravel & 8 & Create, read, update, delete, request validation, flash message, dan handling error berjalan lengkap. & CRUD utama berjalan, tetapi validasi atau handling error belum lengkap. & CRUD tidak berjalan utuh atau data tidak tersimpan benar. \\ \\
Kualitas arsitektur Laravel & 7 & Route resource, controller, model, migration, seeder, dan Blade dipisahkan rapi mengikuti konvensi Laravel. & Struktur cukup mengikuti Laravel, tetapi masih ada logika bercampur atau konvensi tidak konsisten. & Kode tidak mengikuti struktur Laravel atau sulit dipelihara. \\ \\
Pengujian praktik dan defense & 5 & Kasus CRUD normal/invalid diuji, data diverifikasi di database, dan keputusan teknis dapat dijelaskan. & Pengujian dilakukan terbatas atau defense belum menjelaskan semua keputusan. & Tidak ada pengujian bermakna atau tidak mampu menjelaskan modul. \\}

\assessmentblock{Tugas Auth, Import Export, API, dan Testing}{Hands-on Practice}{SCPMK0603-03402, SCPMK1006-03403}{Mahasiswa menyelesaikan tugas auth, import export, api, dan testing sesuai Sub-CPMK dan konteks mata kuliah.}{Menganalisis kebutuhan, mengerjakan artefak praktik/proyek, menguji hasil, mendokumentasikan evidence, dan mempresentasikan keputusan bila diminta.}{Laporan, artefak digital, repository/prototype, evidence pengujian, dan/atau bahan presentasi.}{Autentikasi dan otorisasi & 10 & Login, logout, role/permission, proteksi route, dan akses data diterapkan sesuai skenario. & Auth berjalan, tetapi role, proteksi route, atau akses data masih memiliki celah. & Auth/otorisasi tidak berjalan atau route penting tidak terlindungi. \\ \\
Import, export, dan API & 8.8 & Import/export data, endpoint API, request/response, validasi, dan error handling bekerja konsisten. & Fitur import/export/API berjalan sebagian, tetapi validasi atau error handling belum lengkap. & Import/export/API gagal atau tidak sesuai kontrak data. \\ \\
Testing dan dokumentasi evidence & 6.2 & Feature/unit test, dokumentasi endpoint, contoh request, dan evidence eksekusi tersedia jelas. & Testing atau dokumentasi tersedia, tetapi cakupan dan evidence masih terbatas. & Tidak ada testing memadai atau evidence tidak membuktikan fitur. \\}

\assessmentblock{PBL Aplikasi Laravel Terintegrasi}{Project Based Learning}{SCPMK0603-03402, SCPMK1006-03403}{Mahasiswa menyelesaikan pbl aplikasi laravel terintegrasi sesuai Sub-CPMK dan konteks mata kuliah.}{Menganalisis kebutuhan, mengerjakan artefak praktik/proyek, menguji hasil, mendokumentasikan evidence, dan mempresentasikan keputusan bila diminta.}{Laporan, artefak digital, repository/prototype, evidence pengujian, dan/atau bahan presentasi.}{Kesesuaian kebutuhan dan rancangan aplikasi & 4 & Kebutuhan pengguna, model data, role, dan prioritas fitur Laravel dirancang konsisten. & Rancangan cukup sesuai, tetapi sebagian kebutuhan atau role belum terpetakan jelas. & Aplikasi tidak menunjukkan kebutuhan pengguna atau rancangan tidak konsisten. \\ \\
Integrasi fitur Laravel & 3.5 & CRUD, auth, API/import-export, validasi, relasi database, dan Blade/API terintegrasi stabil. & Fitur utama terintegrasi, tetapi masih ada kekurangan validasi, relasi, atau error handling. & Fitur utama tidak terintegrasi atau aplikasi sering gagal. \\ \\
Testing, deployment, dan presentasi proyek & 2.5 & Test, seeder/demo data, deployment atau demo lokal, repository, dan presentasi keputusan tersedia lengkap. & Evidence proyek tersedia, tetapi testing/deployment atau presentasi keputusan belum kuat. & Tidak ada evidence memadai atau proyek tidak dapat didemonstrasikan. \\}

\assessmentblock{UAS Evaluasi Proyek Laravel}{UAS praktik dan defense}{SCPMK1006-03403}{Mahasiswa menyelesaikan uas evaluasi proyek laravel sesuai Sub-CPMK dan konteks mata kuliah.}{Menganalisis kebutuhan, mengerjakan artefak praktik/proyek, menguji hasil, mendokumentasikan evidence, dan mempresentasikan keputusan bila diminta.}{Laporan, artefak digital, repository/prototype, evidence pengujian, dan/atau bahan presentasi.}{Evaluasi kelengkapan fitur akhir & 4 & Fitur Laravel akhir memenuhi scope, role, alur bisnis, validasi, dan pengelolaan data. & Fitur utama memenuhi scope, tetapi beberapa validasi atau role masih belum lengkap. & Fitur akhir tidak memenuhi scope atau banyak alur gagal. \\ \\
Kualitas teknis dan maintainability & 3.5 & Kode mengikuti konvensi Laravel, struktur rapi, query efisien, dan konfigurasi siap dijalankan ulang. & Kode cukup rapi, tetapi masih ada duplikasi, query kurang efisien, atau konfigurasi kurang jelas. & Kode sulit dipelihara atau tidak dapat dijalankan ulang. \\ \\
Defense, evidence, dan rekomendasi & 2.5 & Defense menjelaskan arsitektur, testing, deployment, risiko, dan rekomendasi peningkatan aplikasi. & Defense cukup jelas, tetapi evidence testing/deployment atau rekomendasi masih terbatas. & Tidak mampu mempertahankan keputusan teknis atau evidence tidak memadai. \\}

\begin{tblr}{width=\textwidth,colspec={|Q[l,m,3.2cm]|X[l,m]|},hlines,vlines,hspan=minimal,row{1-3}={bg=headgray},rowsep=1.5pt,colsep=3pt}
Jadwal Pelaksanaan: & Minggu 1--17 sesuai RPS. Setiap evaluasi memiliki RTM dan rubrik multi-kriteria. \\
Lain-Lain Yang Diperlukan: & LMS, repositori/prototype/proyek, perangkat lunak sesuai mata kuliah, dan perangkat presentasi. \\
Pustaka: & Mengacu pustaka utama dan pendukung pada bagian RPS. \\
\end{tblr}
